\documentclass[11pt]{article}
\usepackage{lecturenotes}

% Uncomment to change the colour of comments in code for this document only.
% \definecolor{codecomment}{gray}{0.25}

\course{CPSC 221: Basic Algorithms and Data Structures}

\begin{document}

\lecture{Insertion Sort}{September 25, 2026}

In the in-class lectures and Friday video, we skipped the correctness proof of insertion sort. We
thought it would be most useful to see a complete proof written out, instead of explained in a
lecture or video. Nathan and Cinda will each write the proof in their own way, so that you can see
two different ways of writing a proof.


Insertion sort is the following algorithm:

\begin{cpp}
// Sorts C
void insertion_sort(vector<int>& C){
    for(i=0; i < C.size()-1; i++){
        insert(C, i+1).
    }
}
\end{cpp}

The \cppinline{insert} subroutine is implemented as follows:

\begin{cpp}
// The following function assumes that:
// 1. loc is in {0, ..., n-1}  where n is the number of entries in C; and
// 2. C[0,...,loc-1] is sorted.
void insert(vector<int>& C, int loc){
    int j = loc;
    while(j > 0 && C[j-1] > C[j]){
        swap(C, j-1, j); // int temp = C[j]; C[j] = C[j-1]; C[j-1] = temp;
        j--;
    }
}
\end{cpp}

\section{Nathan's Proof}

We will prove the theorem in excruciating detail. It is usually not necessary to go into such
excruciating detail for simple algorithms like this, but we will do it because suffering can be good
for you.



\subsection{The \texttt{insert} Subroutine}
We will first analyze the \cppinline{insert} subroutine and then apply it to prove the correctness
of insertion sort.

\newcommand{\cv}[1]{ \ifmmode{\text{\cppinline{#1}}}\else{\cppinline{#1}} \fi }

\begin{lemma}
\label{lemma:nh}
Let $\cv{C}$ be any array of integers and let $\cv{loc} \in \{0, \dotsc, n-1\}$ where
$n = \cv{C.size()}$. Assume that $\cv{C[0, ..., loc-1]}$ is sorted. Then, upon termination
of $\cv{insert(C, loc)}$, $C[0, ..., loc]$ is sorted.
\end{lemma}
\begin{proof}
We first argue that the algorithm will terminate. Since \cv{j} is initialized to \cv{loc}, is
decremented in line 8, and never changed otherwise, the loop will run at most $\cv{loc}$ number of
times, since it requires $\cv j > 0$. 

Let $C_0$ be the initial state of the variable $\cv C$. We define the following loop invariant:
\[
\text{\textbf{LI:}} \begin{cases}
1. &\text{ $\cv C$ is a permutation of $C_0$;} \\
2. &\text{ $\cv{C[j, ..., loc]}$ is sorted;} \\
3. &\text{ If $0 < \cv j \leq \cv{loc}$ then $C[0, ..., j-1]$ is sorted; } \\
4. &\text{ If $0 < \cv j < \cv{loc}$ then $\cv{C[j-1] < C[j+1]}$.}
\end{cases}
\]
Suppose that \textbf{LI} holds at the beginning of the final execution of line 6.  Then line 6
guarantees that, on termination, either:
\begin{enumerate}
\item $j = 0$. In this case, item 2 of \textbf{LI} guarantees that \cv{C[0, ..., loc]} is sorted,
as desired.
\item $j > 0$ but $\cv{C[j-1]} \leq \cv{C[j]}$. In this case, item 2 of \textbf{LI} guarantees that
\cv{C[0, ..., j-1]} and \cv{C[j, ... loc]} are sorted, and $\cv{C[j-1]} \leq \cv{C[j]}$. Therefore
$\cv{C}$ is sorted, as desired.
\end{enumerate}
Now we prove the following claim:
\begin{claim}
Statement \textbf{LI} is true at the start of line 6, every time line 6 is executed.
\end{claim}
\begin{proof}[Proof of claim]
\textbf{Base case:} We first prove that \textbf{LI} holds when line 6 is first executed.
We verify the properties of \textbf{LI}:
\begin{enumerate}
\item $\cv C = C_0$ so it is the trivial permutation.
\item Since $\cv j = \cv{loc}$, this statement trivially holds.
\item $\cv{C[0, ..., loc-1]}$ is sorted by assumption.
\item Since $\cv j = \cv{loc}$, this statement trivially holds.
\end{enumerate}

\textbf{Inductive step:} We show that, assuming \textbf{LI} holds at the start of an execution of
the loop (i.e. after line 6 is executed but before line 7 is executed), then it also holds at the
end of the loop (i.e. after line 8 is executed). This ensures that it is true the next time line 6
starts.

For notational convenience, let $C$ be the state of \cv C at the start of the loop, and let $C'$ be
the state of \cv C at the end of the loop. Let $j$ be the value of \cv j at the start of the loop
and $j' = j-1$ the value of \cv j at the end of the loop. By induction, we assume \textbf{LI} holds
for $C$ and $j$, and we must verify \textbf{LI} for the values $C'$ and $j'$:
\begin{enumerate}
\item $C'$ is a permutation of $C_0$ because $C$ is a permutation of $C_0$, and $C'$ is obtained
from $C$ by swapping two entries.

\item We want to show that $C'[j', \dotsc, \cv{loc}] = C'[j-1, j, \dotsc, \cv{loc}]$ is sorted.
Since entries $j-1, j$ of $C$ were swapped to obtain $C'$,
\[
C'[j-1, j, \dotsc, \cv{loc}]
= \begin{cases}
    (C[j], C[j-1]) &\text{ if } j = \cv{loc} \\
    (C[j], C[j-1], C[j+1], \dotsc, C[\cv{loc}]) &\text{ if } j < \cv{loc} .
  \end{cases}
\]
since entries $j-1, j$ were swapped to obtain $C'$. By the loop condition, $C[j] < C[j-1]$. If $j =
\cv{loc}$ then we are done. Otherwise, $0 < j < \cv{loc}$ so $C[j, \dotsc, \cv{loc}]$ is sorted and
$C[j-1] \leq C[j+1]$ (using \textbf{LI} for $C, j$). Therefore $C[j] < C[j-1] \leq C[j+1] \leq
\dotsm \leq C[\cv{loc}]$, as desired.

\item We want to show that, if $0 < j' \leq \cv{loc}$, then  $C'[0, \dotsc, j'-1]$ is sorted. But
$C'[0, \dotsc, j'-1] = C'[0, \dotsc, j-2] = C[0, \dotsc, j-2]$ since this part of the array is not
modified. So the claim holds by \text{LI} on $C, j$.

\item We want to show that $C'[j'-1] < C'[j'+1]$ if $0 < j' < \cv{loc}$. We have
\begin{align*}
    C'[j'-1]
    &= C'[j-2] \\
    &= C[j-2] \tag{Since this entry is untouched} \\
    &\leq C[j-1] \tag{Since $C[0, \dotsc, j-1]$ is sorted} \\
    &= C'[j] \tag{Since this is the swapped value} \\
    &= C'[j'+1] .
\end{align*}
\end{enumerate}
This concludes the proof of the claim.
\end{proof}
The claim guarantees that \textbf{LI} holds at the beginning of the final execution of line 6, so
the output is correct.
\end{proof}

\subsection{The Whole Algorithm}

\begin{theorem}
Let \cv C be an integer array. Upon termination of \cv{insertion_sort(C)}, \cv C is sorted.
\end{theorem}
\begin{proof}
We define another loop invariant:
\[
    \textbf{LI: } \qquad \cv{C[0, ..., i]} \text{ is sorted.}
\]
We want to establish the following claim:
\begin{claim}
At the beginning of each execution of line 3, \cv{insertion_sort}(C), \textbf{LI} is true.
\end{claim}
\begin{proof}[Proof of claim]
Let $n = \cv{C.size()}$ be the number of elements in array \cv C.

\textbf{Base case:} We want to show that \textbf{LI} is true on the first execution of line 3.
In this case $\cv i = 0$ so $\cv{C[0, ..., i]} = \cv{C[0]}$ so the statement is trivial.

\textbf{Inductive step:} Assume \textbf{LI} is true at the beginning of a loop. We want to show that
it is also true at the end of the loop. Let $C, i$ be the values of \cv C, \cv i at the start of the
loop and let $C', i'$ be the values at the end of the loop (so $i' = i + 1$). By the inductive
hypothesis, $C[0, \dotsc, i]$ is sorted. We want to show that $C'[0, \dotsc, i+1]$ is sorted.  Since
$C[0, \dotsc, i]$ is sorted, the assumption of \cref{lemma:nh} is satisfied. Since $C'[0, \dotsc,
i+1]$ is the value of \cv C after \cv{insert(C,i)} is run on values $C, i$, $C'[0, \dotsc, i+1]$ is
sorted.
\end{proof}

Since \cv i gets incremented in each loop and the loop terminates when $\cv i \geq n$, the algorithm
will terminate. On the final execution of line 3, $\cv{i} = n-1$ and \textbf{LI} is true, so
\cv{C[0, ..., n-1]} is sorted, as desired.
\end{proof}

\end{document}
